\begin{prob}
    % Write the problem below. Precede any paths with \parentdir. For example:
    %   \includegraphics{\parentdir/public/fig/picture.png}
    % Any files used in the solution should be placed in the `private` 
    % subdirectory. Any other files can go in the `public` subdirectory.

    In this problem, we will show that the average case time complexity of
    binary search is $\Theta(\log n)$.

    \begin{subprobset}
        \begin{subprob}
            The number of calls to \python{binary_search} required to find a target depends
            on the target's position within the array. For instance, if the target is
            in the exact middle of the array, only one call to binary search is necessary.
            If the target is elsewhere, more calls are required.

            The graphic below shows 15 blanks, one for each element of a 15-element array.
            In each blank, write the number of calls needed to \python{binary_search}
            in order to find the target if the target were located at that position in
            the array. Several of the blanks have been filled in for you.


            \newcommand{\arrslot}[2]{
                \tikz{
                    \draw (0,0) -- (.5,0);
                    \node at (.25,-.4) {\small \texttt{#1}};
                    \node at (.25,.4) { #2 };
                }
            }

            \begin{center}
                \arrslot{0}{}
                \arrslot{1}{}
                \arrslot{2}{}
                \arrslot{3}{2}
                \arrslot{4}{}
                \arrslot{5}{}
                \arrslot{6}{}
                \arrslot{7}{1}
                \arrslot{8}{}
                \arrslot{9}{}
                \arrslot{10}{}
                \arrslot{11}{2}
                \arrslot{12}{}
                \arrslot{13}{}
                \arrslot{14}{}
            \end{center}


        \end{subprob}

        \begin{subprob}
            Given an array of size $n$, what is the greatest number of calls to
            \python{binary_search} necessary to find a target in the array? Your
            answer should be a formula involving $n$.
            You should assume that the target is in the array, and that $n =
            2^{k}-1$ for some $k \in \{1, 2, 3, \ldots\}$; that is, $n \in \{1,
            3, 7, 15, \ldots\}$.

            \begin{soln}
                If $n = 1$, only one call is necessary.

                If $n = 3$, two calls might be necessary. Note that this is $\log_2 (3 + 1)$.

                If $n = 7$, three calls might be necessary. Note that this is $\log _2 (7 + 1)$.

                Apparently, $\log_2 (n + 1)$ calls are necessary.
            \end{soln}
        \end{subprob}


        \begin{subprob}
            Given an array of size $n$, where $n = 2^{k} - 1$ for some $k$, let
            $f_n(k)$ be the number of elements of the array which require exactly $k$
            calls to \python{binary_search} before they are found. Derive a
            formula for $f_n(k)$. 

            \begin{soln}
                Only one element requires just one call: the middle element. So
                $f_n(1) = 1$.

                Two elements can be found with exactly two calls, so $f_n(2) = 2$.
                
                Four elements can be found with exactly three class, so $f_n(3) = 4$.

                It looks like the number of elements is doubling with each step.
                Here's a table capturing the pattern.

                \begin{center}
                \begin{tabular}{cc}\toprule
                    $k$ & $f_n(k)$\\\midrule
                    1 & 1\\
                    2 & 2\\
                    3 & 4\\
                    4 & 8\\
                    5 & 16\\
                    $\vdots$ & $\vdots$\\\bottomrule
                \end{tabular}
                \end{center}

                So it looks like $f_n(k) = 2^{k-1}$.
            \end{soln}
        \end{subprob}

        \begin{subprob}
            Define
            \[
                S(n) = \sum_{k=1}^{\log_2 (n + 1)} k \cdot 2^{k-1}.
            \]
            Show that $S(n) = O(n \log n)$.

            Hint 1: show that $S(n)$ is smaller than something that is $\Theta(n \log n)$.
            The following fact might help:
            \[\sum_{k=1}^{\log (n+1)} k \cdot 2^{k-1} \leq \sum_{k=1}^{\log (n+1)} \log(n+1) \cdot 2^{k-1}\]

            Hint 2: 
            You may need to remember the formula for the sum of a geometric progression
            from calculus:
            \[
                \sum_{k=0}^K x^k = \frac{1 - x^{K+1}}{1 - x}
            \]

            \begin{soln}
                \begin{align*}
                    S(n) 
                        &= \sum_{k=1}^{\log_2 (n + 1)} k \cdot 2^{k-1}
                    \intertext{We'll try to get an upper bound using hint 1:}
                        \\
                        &\leq
                            \sum_{k=1}^{\log_2 (n+1)} \log_2 (n+1) \cdot 2^{k-1}
                            \\
                        &=
                            \log_2(n+1) \sum_{k=1}^{\log_2 (n+1)} 2^{k-1}
                            \\
                    \intertext{We recognize the sum of a geometric progression and want to
                        use the formula in hint 2 above, but this sum starts at $k=1$
                        instead of $k=0$. We can start the sum at $k=0$ be reindexing:
                    }
                        &=
                            \log_2(n+1) \sum_{k=0}^{\log_2 (n+1) - 1} 2^{k}
                            \\
                    \intertext{Using the formula for the sum of a geometric progression:}
                        &=
                            \log_2(n+1) \left[ \frac{1-2^{\log_2(n+1)}}{1 - 2} \right]
                            \\
                        &=
                            \log_2(n+1) \left(2^{\log_2(n+1)} - 1\right)
                            \\
                    \intertext{You can convince yourself that $\log_2(n+1) = \Theta(\log n)$
                        and $2^{\log_2(n+1)} = \Theta(2^{\log_2 n}) = \Theta(n)$ using
                        any of the techniques that we saw when learning about asymptotic
                        notation, such as limits.}
                        &=
                            \Theta(\log n) \cdot \Theta(n)
                            \\
                        &=
                            \Theta(n \log n)
                \end{align*}
            \end{soln}
        \end{subprob}


        \begin{subprob}
            Define
            \[
                S(n) = \sum_{k=1}^{\log (n+1)} k \cdot 2^{k-1}.
            \]
            Show that $S(n) = \Omega(n \log n)$.

            \begin{soln}
                If we drop all but the last term from the sum, we make it smaller. This
                gives us the lower bound we need:
                \begin{align*}
                    S(n)
                    &=
                        \sum_{k=1}^{\log (n+1)} k \cdot 2^{k-1}.
                        \\
                    &\geq
                        \log(n+1) \cdot 2^{\log_2(n+1)}
                        \\
                    &=
                        \Theta(\log n) \cdot \Theta(n)
                        \\
                    &=
                        \Theta(n \log n)
                \end{align*}
            \end{soln}
        \end{subprob}

        \begin{subprob}
            What is the average case time complexity of binary search? Assume that the
            target is in the array exactly once and that the probability of it being in
            any given position is simply $1/n$.
            shuffled randomly before the first call. You may also assume that the size
            of the array is $2^k - 1$ for some $k \in {1,2,3,\ldots}$.

            Hint: since the time taken during any call to
            \python{binary_search}, excluding time spent in recursive calls, is constant,
            the total time taken by binary search is proportional to the number of calls
            made.

            \begin{soln}
                The average case time complexity can be found by computing the average
                number of calls to \python{binary_search} needed, assuming that the target
                is placed uniformly at random within the array.

                The first case is that only one call is necessary. There is only 
                $f_n(1) = 1$ place where the target can be for this to happen, and this
                occurs with probability $1/n$.

                The second case is that exactly two calls are necessary. There are
                $f_n(2) = 2$ places where the target can be for this to happen,
                so this occurs with probability $2/n$.

                The third case is that exactly three calls are necessary. There are
                $f_n(3) = 4$ places where the target can be for this to happen, so
                this occurs with probability $4/n$.

                In general, case $k$ is when exactly $k$ calls are necessary. There
                are $f_n(k) = 2^{k-1}$ places where the target can be for this to
                happen, so this occurs with probability $2^{k-1} / n$.

                In total, there are $\log_2(n + 1)$ cases.

                The average number of calls needed is therefore:
                \begin{align*}
                    \sum_{k=1}^{\log_2(n+1)} k \cdot \frac{f_n(k)}{n}
                    &=
                        \frac{1}{n}\sum_{k=1}^{\log_2(n+1)} k \cdot 2^{k-1}
                        \intertext{From the previous parts, we know that the
                        summation is both $O(n \log n)$ and $\Omega(n \log n)$,
                        and so it is $\Theta(n \log n)$:}
                    &=
                        \frac{1}{n} \cdot \Theta(n \log n)
                        \\
                    &=
                        \Theta(\log n)
                \end{align*}
                This shows that the average case time complexity is $\Theta(\log n)$.
            \end{soln}
        \end{subprob}
    \end{subprobset}

\end{prob}
